% 出版级排版设置
\usepackage{xeCJK}
\usepackage{geometry}
\usepackage{titlesec}
\usepackage{titletoc}
\usepackage{fancyhdr}
\usepackage{setspace}
\usepackage{indentfirst}
\usepackage{etoolbox}

% 中文目录标题
\renewcommand{\contentsname}{目\quad 录}

% 页边距：装订侧稍宽
\geometry{
  top=2.8cm,
  bottom=2.8cm,
  left=3.2cm,
  right=2.8cm,
  headheight=1cm,
  headsep=0.8cm,
  footskip=1cm
}

% 字体设置
\setmainfont{DejaVu Serif}
\setsansfont{DejaVu Sans}
\setCJKmainfont{Noto Serif CJK SC}[
  BoldFont=Noto Serif CJK SC Bold,
  ItalicFont=Noto Serif CJK SC
]
\setCJKsansfont{Noto Sans CJK SC}[
  BoldFont=Noto Sans CJK SC Bold
]
\setCJKmonofont{Noto Sans Mono CJK SC}
% 八卦符号和带圈数字fallback
\usepackage{newunicodechar}
\newfontfamily{\symbolfont}{DejaVu Sans}
\newfontfamily{\symbolfontbold}{DejaVu Sans Bold}
\newunicodechar{①}{{\symbolfont ①}}
\newunicodechar{②}{{\symbolfont ②}}
\newunicodechar{③}{{\symbolfont ③}}
\newunicodechar{䷗}{{\symbolfont ䷗}}
\newunicodechar{☰}{{\symbolfont ☰}}
\newunicodechar{☱}{{\symbolfont ☱}}
\newunicodechar{☲}{{\symbolfont ☲}}
\newunicodechar{☳}{{\symbolfont ☳}}
\newunicodechar{☴}{{\symbolfont ☴}}
\newunicodechar{☵}{{\symbolfont ☵}}
\newunicodechar{☶}{{\symbolfont ☶}}
\newunicodechar{☷}{{\symbolfont ☷}}

% 行距：1.5倍
\onehalfspacing

% 段落首行缩进2字符
\setlength{\parindent}{2em}
\setlength{\parskip}{0.3em}

% ========== 封面页 ==========
\renewcommand{\maketitle}{%
  \begin{titlepage}
    \centering
    \vspace*{3cm}
    {\sffamily\bfseries\fontsize{42pt}{50pt}\selectfont 生命论\par}
    \vspace{0.5cm}
    {\sffamily\bfseries\fontsize{28pt}{36pt}\selectfont （明本论）\par}
    \vspace{2cm}
    {\sffamily\fontsize{18pt}{24pt}\selectfont ——从操作出发的存在论革命\par}
    {\sffamily\fontsize{18pt}{24pt}\selectfont 与旧哲学总清算\par}
    \vspace{3cm}
    {\sffamily\fontsize{16pt}{22pt}\selectfont 全本·全九卷 + 附录六种\par}
    \vspace{1cm}
    {\sffamily\fontsize{14pt}{20pt}\selectfont 2026年8月\par}
    \vfill
  \end{titlepage}
}

% ========== 卷标题格式（# 第X卷）==========
\titleformat{\part}[display]
  {\centering\sffamily\bfseries}
  {\fontsize{36pt}{44pt}\selectfont\thepart}
  {1em}
  {\fontsize{28pt}{36pt}\selectfont}
\titlespacing*{\part}{0pt}{6cm}{2cm}

% ========== 篇标题格式（## 第X篇）==========
\titleformat{\section}
  {\centering\sffamily\bfseries\Large}
  {\thesection}{1em}{}
\titlespacing*{\section}{0pt}{2em}{1.5em}

% ========== 章标题格式（### 第X章）==========
\titleformat{\subsection}
  {\centering\sffamily\bfseries\large}
  {\thesubsection}{1em}{}
\titlespacing*{\subsection}{0pt}{1.5em}{1em}

% 让每卷另起一页
\pretocmd{\part}{\cleardoublepage}{}{}
\pretocmd{\section}{\clearpage}{}{}

% ========== 页眉页脚 ==========
\pagestyle{fancy}
\fancyhf{}
\fancyhead[LE]{\small\sffamily 生命论（明本论）}
\fancyhead[RO]{\small\sffamily\leftmark}
\fancyfoot[C]{\thepage}
\renewcommand{\headrulewidth}{0.4pt}
\renewcommand{\footrulewidth}{0pt}

% 卷页、章页不显示页眉
\fancypagestyle{plain}{
  \fancyhf{}
  \fancyfoot[C]{\thepage}
  \renewcommand{\headrulewidth}{0pt}
}

% ========== 目录格式 ==========
\setcounter{tocdepth}{1}
\titlecontents{part}
  [0em]
  {\vspace{1em}\bfseries\large}
  {第\thecontentslabel 卷\quad}
  {}
  {\hfill\contentspage}
\titlecontents{section}
  [2em]
  {\vspace{0.3em}}
  {第\thecontentslabel 篇\quad}
  {}
  {\titlerule*[0.5pc]{.}\contentspage}
\titlecontents{subsection}
  [4em]
  {\vspace{0.2em}\small}
  {第\thecontentslabel 章\quad}
  {}
  {\titlerule*[0.5pc]{.}\contentspage}

% 链接颜色
\usepackage{hyperref}
\hypersetup{
  colorlinks=true,
  linkcolor=black,
  urlcolor=black,
  citecolor=black,
  bookmarksnumbered=true
}

% 列表间距
\usepackage{enumitem}
\setlist{nosep, leftmargin=2em}

% 表格
\usepackage{longtable}
\usepackage{booktabs}
\usepackage{array}

% 禁止孤行寡行
\widowpenalty=10000
\clubpenalty=10000

% 强制目录只到篇级
\AtBeginDocument{\setcounter{tocdepth}{1}}
